% SHARD-ID: RC-02D-DEFINITIONS-COMPLEXES
% SHARD-TITLE: Definitions, Continued: Ordered and Pointed Geodesic Flows, Buildings, Twisted Complexes
% SHARD-SUMMARY: Fourth definitions shard (split for length): the two competing higher-dimensional non-backtracking flows -- the ordered-cell flow of an arbitrary complex and the pointed opposition flow of a building -- the graded geodesic determinant that grades them by cell dimension, and the building vocabulary (colour-shift Hecke operators, Ramanujan complex, Cayley complex).
% SHARD-SUMMARY: Then the twisting data: an arbitrary Kraus connection on the 1-skeleton with its flatness and pairing conditions, the voltage quotient and its Artin block, and the quantum vertex property of type A-tilde with its exceptional space.
% SHARD-SUMMARY: Every row of db/definitions.tsv still appears exactly once across the four definitions shards.
% SHARD-KEYWORDS: definitions, ordered cell, geodesic flow, pointed facet, opposition, algebraic length, graded geodesic determinant, superdeterminant, Bruhat-Tits building, Hecke operator, Ramanujan complex, Cayley complex, voltage, Artin block, quantum expander
% SHARD-NOTES: none
% SHARD-SCRIPTS: none

\section{Definitions, Continued: Flows on Complexes and their Twists}
\label{sec:definitions-d}

Continuation of \cref{sec:definitions-c}; split for length only. Throughout,
$X$ is a finite abstract simplicial complex of dimension $d$ with face numbers
$f_0,\dots,f_d$ and Euler characteristic $\chi(X) = \sum_i(-1)^if_i$, and
$s_k = (-1)^{k+1}$.

\begin{definition}[Ordered $k$-cell, ordered geodesic flow]
\label{def:ordered-geodesic-flow}
An \emph{ordered $k$-cell} of $X$ is a tuple $[v_0,\dots,v_k]$ of distinct
vertices whose underlying set is a $k$-simplex of $X$. Write $\mathcal O_k$ for
the set of ordered $k$-cells, so $|\mathcal O_k| = (k+1)!\,f_k$, and
$H_k = \mathbb C^{\mathcal O_k}$. The \emph{ordered geodesic flow} is the
operator $T_k$ on $H_k$, $1\le k\le d$,
\[
 T_k[v_0,\dots,v_k] \;=\; \sum_{\substack{w\notin\{v_0,\dots,v_k\}\\
 \{v_1,\dots,v_k,w\}\in X,\;\{v_0,\dots,v_k,w\}\notin X}} [v_1,\dots,v_k,w] :
\]
drop the oldest vertex, append a new one keeping the ordered overlap, and forbid
the step that a $(k+1)$-cell would short-cut. For $d=1$ this is exactly the
Hashimoto operator of \cref{def:hashimoto-operator}, and the excluded vertex is
the immediately previous one.
\end{definition}

\begin{definition}[Pointed facet, opposition successor, algebraic length]
\label{def:pointed-opposition-flow}
Let $X$ be a finite quotient of the Bruhat--Tits building $\mathcal B$ of
$\mathrm{PGL}_n$ over a local field with residue order $\qs$, carrying the
cyclic type structure of the affine diagram $\tilde A_{n-1}$. A \emph{pointed
$k$-facet} is a $k$-simplex with a distinguished vertex; ordering it cyclically
in increasing type from that vertex gives its \emph{ordered type}, a positive
composition $(\lambda_0,\dots,\lambda_k)$ of $n$, and its \emph{algebraic step
length} is $\lambda_0$. Write $P_k$ for the set of pointed $k$-facets, so
$|P_k| = (k+1)f_k$. A \emph{successor} shifts the ordering and crosses to a
vertex \emph{opposite} in the spherical link of the common $(k-1)$-facet; the
\emph{pointed opposition flow} is
\[
 \mathsf T_k(\uz)[\sigma] \;=\; \uz^{\lambda_0(\sigma)} % bare-ok
 \sum_{\tau\in\operatorname{Succ}(\sigma)}[\tau] \qquad\text{on }\mathbb C^{P_k}. % bare-ok
\]
Opposition is strictly stronger than the non-cell condition of
\cref{def:ordered-geodesic-flow}, and the exponent $\uz^{\lambda_0}$ is a
\emph{length}, never a matrix power or a multiplicity.
\end{definition}

\begin{definition}[Graded geodesic determinant, total zeta of a complex]
\label{def:graded-geodesic-determinant}
\emph{Geodesic data} $\mathfrak g$ on $X$ is, in each degree $1\le k\le d$: a
finite state set $\mathcal S_k$, a successor relation with incidence
multiplicities, a positive integer algebraic length $\wl_k(\sigma)$, and a step % bare-ok
transport $U_{\sigma\tau}$ on a finite-dimensional fibre $V$. Give the whole % bare-ok
degree-$k$ space $V\otimes\mathbb C^{\mathcal S_k}$ parity $k+1$, put
$\mathsf T^U_k(\uz)(v\otimes[\sigma]) = \sum_{\sigma\rightsquigarrow\tau} % bare-ok
\uz^{\wl_k(\sigma)}U_{\sigma\tau}v\otimes[\tau]$, and let % bare-ok
$W = \bigoplus_{k=1}^d V\otimes\mathbb C^{\mathcal S_k}$. The \emph{graded
geodesic determinant} (the total zeta) is
\[
 \gzeta(X,\mathfrak g,U;\uz)
 \;=\; \sdet_W\Bigl(I - \bigoplus_{k=1}^d s_k\mathsf T^U_k(\uz)\Bigr)^{-1}
 \;=\; \prod_{k=1}^d\det\bigl(I - s_k\mathsf T^U_k(\uz)\bigr)^{(-1)^k},
\]
with the sign $s_k = (-1)^{k+1}$ inserted once per successor arrow. Odd $k$
(edges) sits in the denominator, even $k$ (chambers) in the numerator. The
intrinsic default for an arbitrary complex is $\mathcal S_k = \mathcal O_k$,
unit length, the rule of \cref{def:ordered-geodesic-flow} and $V=\mathbb C$,
written $\gzeta_{\mathrm{ord}}(X,\uz)$; the building default is
\cref{def:pointed-opposition-flow}, written $\gzeta_{\mathrm{pt}}(X,\uz)$.
\end{definition}

\begin{definition}[Bruhat--Tits building of $\mathrm{PGL}_n$, colour-shift Hecke operators, Ramanujan complex]
\label{def:bruhat-tits-building}
The vertices of $\mathcal B$ carry a colour (type) function with values in
$\mathbb Z/n\mathbb Z$, and the \emph{colour-shift Hecke operators} $A_k$,
$1\le k\le n-1$, sum a function over the neighbours of colour shift $k$:
{\small\texttt{\detokenize{we may look at the $d-1$ colored adjacency operators $A_k$, $k = 1,\dots,d-1$ on $\L2{\B^0}$, called the Hecke operators.}}}
(\texttt{math/0406208:main.tex:541-549}). They commute, are normal, and satisfy
$A_k^* = A_{n-k}$. A finite quotient is a \emph{Ramanujan complex} when its
nontrivial joint spectrum lies in the tempered set
$\mathrm{Aspec}_n = \{(\qs^{k(n-k)/2}e_k(z))_k : |z_i| = 1,\ \prod_iz_i = 1\}$:
{\small\texttt{\detokenize{A finite quotient $X$ of $\B$ is called a Ramanujan complex if the eigenvalues of every non-trivial simultanenous eigenvector $v \in \L2{X}$, $A_k v = \lam_k v$, satisfy $(\lam_1,\dots,\lam_{d-1}) \in \Aspec_d$.}}}
(\texttt{math/0406208:main.tex:564-572}). The condition is on the \emph{joint}
spectrum, which is strictly smaller than the product of the separate spectra.
\end{definition}

\begin{definition}[Cayley complex]
\label{def:cayley-complex}
For a finite group $G$ with generating set $S$, the \emph{Cayley complex} is the
clique complex of the Cayley graph:
{\small\texttt{\detokenize{The Cayley complex of $G$ with respect to a set of generators $S$ is the simplicial complex whose $1$-skeleton is the Cayley graph $\Cayley{G}{S}$, where a subset of $i+1$ vertices is an $i$-cell iff every two vertices comprise an edge.}}}
(\texttt{math/0406217:main.tex:497-515}). Lubotzky, Samuels and Vishne give
explicit $S$, graded into colour classes $S_1,\dots,S_{n-1}$ with
$|S_k| = \binom{n}{k}_{\qs}$ and $S_k^{-1} = S_{n-k}$, for which this complex is
a Ramanujan complex covered by $\mathcal B$.
\end{definition}

\begin{definition}[Kraus connection on the 1-skeleton; flat, pure gauge, paired]
\label{def:kraus-connection}
A \emph{connection} on the 1-skeleton of $X$ with fibre $V$ assigns to each
directed edge $(x,y)$ an arbitrary $E_{(x,y)}\in\operatorname{End}(V)$; it is a
\emph{Kraus connection} when $V = \Mn$ and each $E_e$ is completely positive,
for instance $E_e = \Ad(B_e)$. Twist \cref{def:ordered-geodesic-flow} by
anchoring the fibre at the \emph{last} vertex,
$T^E_k(v\otimes[\sigma]) = \sum E_{(v_k,w)}v\otimes[\tau]$. The connection is % bare-ok
\emph{flat} when $E_{yx} = E_{xy}^{-1}$ and $E_{yz}E_{xy} = E_{xz}$ on every
$2$-cell, and \emph{pure gauge} when $E_{(x,y)} = F_yF_x^{-1}$ for some
invertible $F$. It is \emph{adjoint-paired} when $E_{\bar e} = E_e^\dagger$ and
\emph{inverse-paired} when $E_{\bar e} = E_e^{-1}$
(\cref{def:kraus-pairings}). No positivity, invertibility, pairing or flatness
is assumed unless said.
\end{definition}

\begin{definition}[Voltage quotient and Artin block]
\label{def:voltage-quotient}
Let a finite group $G$ act on $X$ preserving the successor and length data and
freely on the state set $\mathcal S_k$. Choose orbit representatives $\sigma_a$ % bare-ok
and encode each transition $\sigma_a\to h\sigma_b$ by its \emph{voltage} % bare-ok
$h\in G$; this is the \emph{voltage quotient}, a digraph on the orbits with a
group label on every arrow. For a unitary $\hol$ of $G$, the
\emph{$\hol$-block} $\mathsf T_{k,\hol}(\uz)$ replaces that arrow by
$\hol(h^{-1})$ on the fibre, keeping the length weight. For
$R = \pi\otimes\bar\pi = \bigoplus_\hol a_\hol\hol$ the \emph{Artin
determinant} is $D_{k,R}(\uz) = \prod_\hol
\det(I-s_k\mathsf T_{k,\hol}(\uz))^{a_\hol}$. This is not the same object
as the connection of \cref{def:kraus-connection} on the full cover.
\end{definition}

\begin{definition}[Ramanujan quantum expander of type $\tilde A_{n-1}$; exceptional space]
\label{def:ramanujan-quantum-expander-a2}
Let $X = \operatorname{Cay}(G,S)$ have the colour classes of
\cref{def:cayley-complex}, let $\pi\colon G\to U(N)$ be unitary and
$R = \Ad\pi$ on $\Mn$. Put
$\chan_k(\dm) = |S_k|^{-1}\sum_{s\in S_k}\pi(s)\dm\pi(s)^\dagger$ and
$\hat A_k = |S_k|\chan_k = \sum_{s\in S_k}R(s)$. Let $\mathcal F_R$ be the
common fixed space and let $\mathcal E_R\supseteq\mathcal F_R$ be the sum,
inside $R$, of the type-character subrepresentations excluded by the convention
of \cref{def:bruhat-tits-building}. The tuple $(\chan_1,\dots,\chan_{n-1})$ is
\emph{Ramanujan of type $\tilde A_{n-1}$ relative to $\mathcal E_R$} when its
joint spectrum on $\mathcal E_R^{\perp}$ lies in $\mathrm{Aspec}_n$. Only when
$\mathcal E_R = \mathcal F_R$ is this the fixed-point-complement condition; on a
partite complex $\mathcal E_R$ can be strictly larger.
\end{definition}
